\subsection{Module A.5: strict-handoff feasibility}

\paragraph{Exact interface.}
Assume the six open V triangles cover $\partial H$.  The body selector is
\[
 \xi_i\in(1-A_{i+1},B_i),
 \qquad
 a_i:=1-\xi_{i-1}<A_i,
 \qquad
 b_i:=\xi_i<B_i.
\]

\paragraph{Variables and feasible set.}
Put
\[
 \ell_i:=1-A_{i+1},
 \qquad
 u_i:=B_i.
\]
The exact incident open traces are $X_i([0,u_i))$ and
$X_i((\ell_i,1])$.  Boundary coverage is therefore equivalent to the six
strict overlap conditions $\ell_i<u_i$.

\paragraph{Objective.}
Choose one $\xi_i\in(\ell_i,u_i)$ on every edge, while preserving the unique
actual supercritical index when $N_+=1$ and producing at least two selected
supercritical indices when $N_+\ge2$.

\paragraph{Exact result.}

\begin{lemma}[Strict-handoff selector]
\label{lem:app-strict-handoff-selection}
The six selector intervals are nonempty.  Every simultaneous choice produces
$0<a_i<A_i$, $0<b_i<B_i$, and
\[
 X_{i-1}(1-a_i),X_i(b_i)\in U_i,
 \qquad
 a_i^2+a_ib_i+b_i^2<1.
\]
If $N_+=1$, every choice has exactly one index with $a_i+b_i>1$, equal to the
actual supercritical index.  If $N_+\ge2$, a simultaneous choice exists with
at least two indices satisfying $a_i+b_i>1$.
\end{lemma}

\paragraph{Solution.}

\begin{proof}
Because $V_i$ is interior to a unit-diameter triangle while each other
distinguished point is at distance at least one,
\[
 0<A_i<1,
 \qquad
 0<B_i<1.
\]
By \zcref{lem:app-exact-trace-normalizations}, only the two incident V
triangles can cover a relative-interior point of $e_{i,i+1}$.  Their traces
cover the edge and are open at their outer endpoints, so they must overlap:
$\ell_i<u_i$.  Thus each selector interval is nonempty.  Its choices put the
two selected anchors in $U_i$; their mutual squared distance is
$a_i^2+a_ib_i+b_i^2$ and is strictly less than the squared diameter one.

The actual and selected ascent tests are
\begin{equation}
 A_i+B_i>1\iff\ell_{i-1}<u_i,
 \qquad
 a_i+b_i>1\iff\xi_{i-1}<\xi_i.
 \label{eq:app-handoff-ascents}
\end{equation}
If actual role $i$ is nonsupercritical, then
\[
 \xi_i<u_i\le\ell_{i-1}<\xi_{i-1},
\]
so every selected version is strictly subcritical.  Also
\begin{equation}
 \sum_{i=0}^5(a_i+b_i)
 =\sum_{i=0}^5(1-\xi_{i-1}+\xi_i)=6.
 \label{eq:app-handoff-telescoping}
\end{equation}
When $p$ is the sole actual supercritical index, the other five selected
sums are below one, so \eqref{eq:app-handoff-telescoping} forces the $p$th
sum above one.  This proves the $N_+=1$ assertion.

Suppose next that $p,q$ are nonadjacent actual supercritical indices.  For
each $r\in\{p,q\}$ choose
\[
 \ell_{r-1}<z_r<u_r,
\]
and then
\[
 \xi_{r-1}\in(\ell_{r-1},\min\{u_{r-1},z_r\}),
 \qquad
 \xi_r\in(\max\{\ell_r,z_r\},u_r).
\]
The two variable pairs are disjoint and give selected ascents at $p$ and
$q$.  Any three indices on a six-cycle include a nonadjacent pair, so this
settles $N_+\ge3$ as well.

It remains that exactly two actual supercritical indices, $p,p+1$, are
adjacent.  Along the other four nonsupercritical roles the successive
inequalities $\ell_i<u_i\le\ell_{i-1}$ imply
$\ell_{p-1}<\ell_{p+1}$.  Combining this with supercriticality at $p,p+1$
gives
\[
 \max\{\ell_{p-1},\ell_p\}<\min\{u_p,u_{p+1}\}.
\]
Choose $\xi_p$ between these bounds and then choose
$\xi_{p-1}<\xi_p<\xi_{p+1}$ in their respective overlap intervals.
Equation \eqref{eq:app-handoff-ascents} makes both selected roles
supercritical.
\end{proof}

\paragraph{Body consequence.}
This supplies every clause, including endpoint strictness, in
\zcref{prop:strict-handoffs}.
